% Folder 134-2, batch 05 — archivists' pages 81–100.
% Pass: Opus 5.5 (claude-opus-5-5), 2026-10-15 — first pass, unchecked against the pages by a human.
%
% Special rules for folder 134-2 (issues #44 and #57, the latter as decided by the
% user on 2026-10-15) applied: Pursuing Stacks text (printed in Maltsiniotis (ed.),
% A la poursuite des champs, vol. I, SMF 2022) is not transcribed, one \page{} and
% one \note{} per page; his own unprinted texts are transcribed in full; texts by
% others are summarised. No third-party text in this batch.
%
% Contents of the batch:
%  - p. 81 (his p. 65): « NOTES to Chapter I and Appendix », end of note (25) and
%    notes (26)-(32), Pursuing Stacks, with his handwritten corrections and a
%    handwritten footnote to note (31). NOT transcribed (one note).
%  - p. 82: blank but for a pencilled « I »; absent.
%  - p. 83: « APPENDIX to Chapter I: three letters to Larry Breen », his English
%    introduction — a second copy of p. 37 (batch 2), in an EARLIER state of
%    correction (no « subtitles App. 1 to App. 18 », no « The later part of the
%    letter »). Transcribed in full, on the same ground as p. 37 (his, not in the
%    printed vol. I).
%  - p. 84: typed « RECOLTES ET SEMAILLES — SOMMAIRE », parts I-IV, sections 1-24,
%    cancelled by a diagonal stroke. Transcribed in full (his own, per the brief).
%  - p. 85 (his p. 63): second copy of Pursuing Stacks notes (12)-(21) (cf. p. 79,
%    batch 4). NOT transcribed (one note).
%  - p. 86: typed « NOTES (*) » list of Récoltes et semailles, notes 1-16 with section
%    numbers and note numbers added in his hand, cancelled by a vertical stroke.
%    Transcribed in full.
%  - p. 87 (his p. 64): second copy of Pursuing Stacks notes (22)-(25), with his
%    corrections. NOT transcribed (one note).
%  - p. 88: typed page « 38 » of a Récoltes et semailles draft (French), heavily
%    corrected in his hand, cancelled by a diagonal stroke. Transcribed in full per
%    the brief. FLAG for the parent: Récoltes et semailles is itself in print
%    (Gallimard, 2022); this page is a draft state, but the parent may want to
%    decide whether the « printed, hence cited not transcribed » rule extends to it.
%  - p. 89 (his p. 65): a clean copy of p. 81 without his handwritten additions.
%    NOT transcribed (one note).
%  - p. 90: an unrelated verso (Bourbaki seminar abstracts, not his, upside down);
%    absent. p. 91: blank; absent.
%  - pp. 92-100 and on: the second typed copy of the letter to Daniel Quillen
%    (« Dear Daniel », 19.2.1983, continued « 20.2. »), Pursuing Stacks §§ 1-4,
%    numbered by hand 1, 1', 2, 2', [3], 3'(?), 4, 4', 5, with the § numbers
%    « 1. » to « 4. » written in the margin. This copy duplicates pp. 14-36; the
%    rule kept is « imprimé, donc cité et non transcrit »: one note per page,
%    naming the §§, his page number, and where his corrections fall (not
%    transcribed in full). The letter runs on past p. 100.
%
% Letters in this batch: letter to D. Quillen (pp. 92-100, his, printed: noted only).
% No letter to Breen and no third-party letter in pages 81-100.
% The brief expected pp. 81-91 to hold typed notes (21)-(31); in fact they hold notes
% (25)-(32) (p. 81, p. 89), (12)-(21) (p. 85) and (22)-(25) (p. 87), interleaved with
% the Récoltes et semailles pages 84, 86, 88 and the appendix introduction p. 83.
\documentclass[11pt,a4paper]{article}
\input{../preamble/grothendieck}

\folder{134-2}
\batch{5}
\pages{81}{100}
\foldertitle{[Chapitre I : Take off (pages 1 à 65) et table des matières provisoire] : tapuscrits annotés (19/02-22/02/1983), note manuscrite (s.d.), copies de lettre (1975, s.d.).}
\dating{1975-[1983]}
\watermark{Édition de démonstration}

\begin{document}

\section{Notes au chapitre I de \emph{Pursuing Stacks} (fin)}

\page{81}
\note{tapuscrit « NOTES to Chapter I and Appendix » (sa p. 65) : fin de la note (25), notes (26) à (32), texte de \emph{Pursuing Stacks}, non transcrit ici : voir G. Maltsiniotis (éd.), \emph{À la poursuite des champs}, vol.~I, SMF, Documents mathématiques 20, 2022. Corrections de sa main : quelques mots retouchés (note (26), « inaccurate »), une phrase biffée dans la note (31), la note (32) cernée d'un trait ; au pied de la page, une note manuscrite appelée depuis la note (31), où il demande si le contre-exemple des types d'homotopie rationnels subsiste quand les $H^i(X)$ sont de dimension finie sur $\mathbb{Q}$, pour un type d'homotopie « de type fini ».}

\section{Appendice au chapitre I : introduction (second exemplaire)}

\page{83}
\note{second exemplaire de l'introduction de l'appendice (cf. p.~37, lot~2), dans un état de correction antérieur : il n'y figure ni l'ajout sur les sous-titres « App. 1 to App. 18 », ni « The later part of the letter ». Tapuscrit anglais, corrections à l'encre de sa main.}

APPENDIX to Chapter I : three letters to Larry Breen

In this appendix, I am encluding\note{sic} three letters to Larry Breen, dated
5.2, 17.2 and 17-19.7. 1975. These letters are written in French, and the
first two are reproduced textually, whereas the third appears here in English
translation. I am grateful to Ronnie Brown who took the trouble last year to
make such a translation from a hardly legible \struck{handwri} copy of the
handwritten letter to Larry Breen, with the assistance of Larry Breen himself
and J.L. Loday. I am now using his translation rather than the original
letter, which no printer could possibly decipher ! Also, for the second letter
I am using a typed copy made in 1975 by Larry Breen (who presumably had
difficulties too deciphering the handwriting), Thanks are due to him for his
interest and patience (with someone like me, very unknowledgeable in standard
homotopy techniques), \struck{as it} \add{which} appeared \struck{from}
\add{in} the letters I got from him in response and his verbal explanations on
related matters, \add{as well as in} the trouble he took in \struck{copying}
retyping \add{that letter} and sending me copies of my own letters to him, and
allowing me to reproduce them in this volume of Pursuing Stacks. My only
present contribution to this set of letters is adding a few comments (in the
notes (22) to (31)), and correcting some inaccuracies in the English
translation (due mainly to my handwriting …) Also, I skipped the beginning of
the first letter, which does'nt seem to me of general relevance.
\note{« as it » biffé est lu avec doute sous la rature. Le « u » de « Pursuing » est retracé à la main.}

The first two letters are an attempt to explain to Larry Breen (who has a wide
background in algebraic geometry and homological and homotopical algebra) some
of the main points of the programme I had in mind around the notions of
n-categories and n-stacks (which is what I am supposed to be pursuing now in my
present work "Pursuing Stacks"). They were written under the impetus of
\struck{the} \add{new} intuition \struck{which} (new to me at any rate) which
then had just appeared to me, namely that (non strict) n-groupoids should model
(in a suitable sense) n-truncated homotopy types.
\note{les parenthèses autour de « which is what … "Pursuing Stacks" » sont tracées à la main, la virgule qui précédait est biffée. Après « homotopy types. » un double trait vertical marque un alinéa.}

The third letter, \add{written in} answer to a number of questions in Larry
Breen's response to the first two, is of a wider scope. A large part of the
letter outlines (very sketchyly) some main points of a duality program
(including a cohomological formulation of "geometric" local and global
classfield theory), which appears here for the first time in print,
\struck{and} going back to the end of the fifties. It gives also some hints
about the need of a framework of "tame topology" suitable for writing up a
"dévissage theory" of stratified spaces, and for \struck{\ill{}}
\add{working} with étale tubular neighbourhoods, for the common purpose of
coming to grips \struck{of} \add{with} a suitable notion of "fine homotopy
type" of a "tame" topological space or \add{of} a scheme, in terms of the
\struck{ordered set} \add{inverse system} of "indexed homotopy types"
corresponding to equisingular stratifications.
\note{« sketchyly » : ainsi dans la frappe, le « t » retouché à la main. « working » est frappé au-dessus du mot biffé, qui est noirci.}

\section{\emph{Récoltes et semailles}, sommaire et notes ; notes de \emph{Pursuing Stacks} (second exemplaire)}

\page{84}
\note{page dactylographiée annulée d'un grand trait oblique ; quelques traits de plume en marge (crochets devant II et devant les sections 9 à 11).}

RECOLTES ET SEMAILLES

SOMMAIRE

I Travail et découverte

1. L'enfant et le Bon Dieu

2. Erreur et découverte

3. Les inavouables labeurs

4. Infaillibilité (des autres) et mépris (de soi)

II Le rêve et le Rêveur

5. Le rêve interdit

6. Le Rêveur

7. L'héritage de Galois

8. Rêve et démonstration

III Naissance de la crainte

9. L'étranger bienvenu

10. La "Communauté mathématique" : fiction et réalité

11. Rencontre avec Claude Chevalley, -- ou : liberté et bons sentiments

12. Le mérite et le mépris

13. Force et épaisseur

14. Naissance de la crainte

15. Récoltes et semailles

IV Le double Visage

16. Marais et premiers rangs

17. Terry Mirkil

18. Vingt ans de fatuité, ou : l'ami infatigable

19. Le monde sans amour

20. Un monde sans conflit ?

21. Un secret de Polichinelle bien gardé

22. Bourbaki, ou ma grande chance -- et son revers

23. De Profundis

24. Mes adieux, -- ou : les étrangers

\page{85}
\note{second exemplaire du tapuscrit des notes (12) à (21) au chapitre I de \emph{Pursuing Stacks} (sa p. 63 ; cf. p.~79, lot~4), les appels des notes (12), (13), (14), (19), (20), (21) suivis d'un « ! ». Texte de \emph{Pursuing Stacks}, non transcrit ici : voir Maltsiniotis (éd.), \emph{À la poursuite des champs}, vol.~I, SMF 2022.}

\page{86}
\note{page dactylographiée annulée d'un grand trait vertical. Les numéros entre parenthèses de la colonne de droite (numéros des notes) sont ajoutés à la main ; les numéros de section qui les précèdent sont frappés. Le tiret qui sépare le titre des numéros est de l'édition.}

NOTES (*)

1. Mes amis de Survivre et Vivre — 11 (\uncertain{5})

2. Aldo Andreotti, Ionel Bucur — 14 (11)

3. Jésus et les douze apôtres — 25 (19)

4. L'Enfant et le Maître — 26 (23)

5. La peur de jouer — 29 (23'')

6. Les deux frères — 29 (23''')

7. Causes d'un échec — 31 (23\textsuperscript{iv})

8. Consensus déontologique -- et contrôle de l'information — 32 (25)

9. Le "snobisme des jeunes", ou les défenseurs de la pureté — 33 (27)

10. Cent fers dans le feu, ou : rien ne sert de sécher ! — 36 (32)

11. L'étreinte impuissante — 37 (34)

12. La visite, ou le rêve malvenu — 45 (40)

13. Krishnamurti, ou \struck{\ill{}} \add{la} libération devenue entrave — 45 (41)

14. L'arrachement salutaire — 45 (42)

15. Mes orphelins — 50 (46a)

Refus d'un héritage, ou le prix d'une contradiction — 50 (46b)

16. On n'arrête pas le progrès ! — 50 (47)
\note{le premier numéro manuscrit, lu « 5 », a la forme d'un $\sigma$ ; le « 19 » est surchargé ; dans « 46a » le dernier caractère est surchargé. La suite saute de (23) à (23'') : ainsi sur la page.}

(*) Voici une liste des notes les plus longues ou les plus substantielles, par ordre d'apparition dans le texte. Les numéros placés après le nom de chaque note (lequel, je l'espère, le rend aisément reconnaissable) indiquent celui de la section (entre 1 et 50) où la note figure, suivi de celui de la note elle-même. Une des notes, plus longue que la longue section à laquelle elle se rapporte, a été subdivisée après coup en deux parties, auxquelles il m'a paru utile de donner des noms séparés.

\page{87}
\note{second exemplaire du tapuscrit des notes (22) à (25) (sa p. 64 ; cf. p.~80, lot~4), avec corrections de sa main dans la note (22) (définition des catégories de Picard, « Picard category is a groupoid … ») et dans la note (24) (« triangulated », « relevant »). Texte de \emph{Pursuing Stacks}, non transcrit ici : voir Maltsiniotis (éd.), \emph{À la poursuite des champs}, vol.~I, SMF 2022. La note (25) se poursuit p.~81 et p.~89.}

\section{\emph{Récoltes et semailles} : une page de brouillon}

\page{88}
\note{page dactylographiée numérotée « 38 » (le chiffre du milieu noirci), corrigée à l'encre de sa main et annulée d'un grand trait oblique. Les mots biffés à la machine qui ne se lisent pas sous les x sont marqués illisibles.}

au cours de cette année-là que j'ai appris pour la première fois (mais
\struck{non} \add{\uncertain{certes pas}} la dernière\add{ !}) \struck{\ill{} (*)} que c'est considéré comme un manque de
sérieux de démontrer bel et bien des choses délicates et visiblement
fondamentales, \add{que tout le monde utilise et} que les prédécesseurs se sont toujours contentés d'admettre\add{(*)}. Ou de démontrer un résultat qui englobe comme cas particuliers plusieurs résultats \add{profonds} connus (ce qui montre bien évidemment que le résultat soi
disant nouveau ne peut être qu'un cas particulier ou une conséquence facile
des résultats connus). Ou de prendre la peine, \struck{\ill{}} \add{dans \struck{un} l'énoncé d'un} résultat
ou dans la description d'une situation en termes d'une autre, de formuler
avec soin les hypothèses naturelles (signe d'un regrettable bombinage),
plutôt que de se \struck{contenter} borner \struck{aux cas simples} à quelque cas d'espèce au
goût de celui qui émet son opinion. Ainsi j'ai vu faire reproche à
Contou-Carrère de ne pas s'être borné dans sa thèse à se placer sur le
corps des complexes ou tout au moins sur un corps, au lieu d'un schéma
général (tout en lui concédant \add{quand même} la circonstance atténuante que c'était \struck{son patron \ill{}} \add{sûrement sur les instances de son patron qu'il avait dû s'y résoudre).} (Celui qui s'exprimait ainsi était pourtant suffisamment dans le
coup pour savoir que même en se bornant au corps des complexes, les nécessités des \struck{\ill{}} démonstrations forçaient la main pour introduire des schémas de
base généraux …)
\note{la ponctuation et les parenthèses de ce passage sont retouchées à la main : la parenthèse ouverte avant « tout en lui concédant » se ferme à la main après « résoudre », et une nouvelle parenthèse s'ouvre avant « Celui qui ». Sous la ligne « Celui qui s'exprimait ainsi », quelques mots manuscrits biffés, illisibles. Dans « (*) » biffé après « dernière ! », l'appel est déplacé à la fin de la phrase suivante.}

\marginal{signe d'alinéa}
Les égarements d'une certaine mode \add{(aujourd'hui)} vont jusqu'à \struck{\ill{}} \struck{exclure} \add{honnir} non seulement \struck{\ill{}} \add{les} démonstrations en \struck{bonne et due} forme, mais \add{souvent} même
des énoncés \struck{en bonne et en forme} \struck{un énoncé en bonne et due forme (sans même parler des hypothèses les plus naturelles), et} et des définitions en forme : au prix
où est le papier et la longanimité du lecteur, il
ne sera bientôt plus question de s'encombrer d'un luxe aussi coûteux\add{ !} \struck{et \ill{} on se bornera prochainement à} En extrapolant les tendances actuelles, il ne doit pas être difficile de
\struck{calculer} \add{prédire} le moment exact où il ne sera plus question \add{dans une publication} d'expliciter \struck{énoncés ni} définitions \add{ni énoncés}, qu'on se \struck{contentera de désigner} \add{désormais de nommer} par des mots-code, en
laissant à l'infatigable et génial lecteur le soin de remplir les blancs,
conformément à ses propres lumières. La tâche du referee sera facilitée
d'autant, car il lui suffira de \struck{consulter} \add{regarder} dans l'annuaire "Who is Who" si
l'auteur est connu comme crédible (de toutes façons personne ne pourrait
contredire les blancs et les \struck{petits points} \add{pointillés} qui composent le brillant
article), ou au contraire un inavouable inconnu qui sera (comme c'est le
cas déjà dès aujourd'hui) \struck{éliminé} \add{éjecté} d'office …
\note{« profonds », « quand même », « honnir », « dans une publication » et « désormais de nommer » sont frappés à la machine dans l'interligne ; les autres ajouts sont à l'encre. « honnir » surmonte le mot noirci après « jusqu'à » ; « exclure », frappé dans l'interligne inférieur, est lui aussi biffé. Le « Les » noirci devant « non seulement » est illisible. « inconnu qui » : un trait sépare les deux mots collés par la frappe.}

(*) Voir aussi à ce sujet l'épisode "La note -- ou la nouvelle éthique"
\add{(par.~33)}. \struck{Qui} \add{(cette fameuse "note")} avait justement le tort d'expliciter des notions \struck{\ill{}} \add{et} des \struck{dém} énoncés \struck{et des démonstrations} qui avaient été laissées jusque là dans le
vague, et qui pourtant ont été \struck{utilisés (par moi \ill{} autant par \ill{} élèves) sans vergogne} \add{implicitement utilisés par moi} pour établir des résultats qui portent mon nom et \struck{\ill{}}
que tout le monde \struck{(y compris \ill{} élève)} utilise sans vergogne\add{,}
\add{depuis bientôt vingt-cinq ans.}
\marginal{sans interligne}
\note{« implicitement utilisés par moi » est frappé dans l'interligne au-dessus de la ligne biffée. La note marginale, cerclée à gauche de la note (*), est une consigne de frappe. Au-dessus du passage biffé « (y compris … élève) » un petit mot manuscrit, lu « mon » avec doute, est lui aussi biffé.}

\section{Notes au chapitre I de \emph{Pursuing Stacks} (exemplaire sans corrections)}

\page{89}
\note{autre exemplaire de la p.~81 (sa p. 65 : fin de la note (25), notes (26) à (32)), sans les corrections ni la note manuscrite de la p.~81. Texte de \emph{Pursuing Stacks}, non transcrit ici : voir Maltsiniotis (éd.), \emph{À la poursuite des champs}, vol.~I, SMF 2022.}

\section{Lettre à Daniel Quillen, 19.2.1983 (second exemplaire) — \emph{Pursuing Stacks}}

\page{92}
\note{second exemplaire dactylographié de la lettre à Daniel Quillen (« Dear Daniel »), qui double les pp.~14 à 36 ; règle suivie : imprimé, donc cité et non transcrit (Maltsiniotis (éd.), \emph{À la poursuite des champs}, vol.~I, SMF 2022). Sa p. « 1 » : § 1, numéro « 1. » ajouté en marge. Ajouts manuscrits dans l'interligne (« program for a », « anyhow », « psychological », une parenthèse sous « algebra ») et quelques ponctuations.}

\page{93}
\note{sa p. « 1' » : fin du § 1 et § 2 (numéro « 2. » en marge) ; quelques corrections de lettres et une surcharge (« who »).}

\page{94}
\note{sa p. « 2 » : § 2 (suite). Ajouts manuscrits : « or rather homotopy » dans l'interligne ; en marge droite, une remarque sur les notations (« NB we write also $\mathrm{id}_a$ instead of $k_1^i(a)$ and $\check{u}$ for $\mathrm{inv}_i(u)$ ») ; au pied, la seconde formule pour $t_1^i(v \ast_\ell u)$, avec « $\ell \geqslant 2$ i.e. $\ell - 1 \geqslant 1$ » ; retouches d'indices dans les formules frappées.}

\page{95}
\note{sa p. « 2' » : § 2 (suite). Les hypothèses de la relation de Godement sont écrites à la main (conditions sur $t_\alpha$, $s_\alpha$, $t_\nu$, $s_\nu$) ; ajouts « previous », « (i-1) » et retouches d'indices.}

\page{96}
\note{sa p. [3] (numéro en tête rogné) : fin du § 2 et § 3 (numéro « 3. » en marge) ; « ever » remplace un mot biffé (« increasingly ») ; la dernière phrase s'achève à la main (« playing the role of the unit segment »).}

\page{97}
\note{sa p. « 3' » (lecture incertaine : le chiffre a la forme d'un 4) : fin de la lettre du 19.2.1983, « Very cordially yours », signée à la main « Alexander ».}

\page{98}
\note{sa p. « 4 », datée « 20.2. » : reprise de la lettre et § 4 (numéro « 4. » en marge) ; ajouts manuscrits « in order » et « $\Pi_\infty$ » au-dessus d'un mot biffé, ponctuation retouchée.}

\page{99}
\note{sa p. « 4' » : § 4 (suite) ; ajouts manuscrits « as non-primitive », « the real », « if viewed as primitive ».}

\page{100}
\note{sa p. « 5 » : § 4 (suite) ; « structures » remplace « choices » biffé, et une note manuscrite en marge gauche, écrite en travers, est appelée dans la phrase sur les comodules (elle commence « or the structure of a » et mentionne un torseur sous un groupe ; lecture partielle). La lettre se poursuit au-delà de ce lot.}

\end{document}
